%% Uses the local my_memo document class (docs/my_memo.cls), which auto-loads
%% graphicx, amsmath, amssymb, hyperref, xcolor, kotex, fancyhdr, and sets its
%% own A4 geometry / boxed title / header-footer.  booktabs and longtable are
%% not provided by the class, so they are loaded here.
\documentclass{my_memo}
\usepackage{booktabs}
\usepackage{array}
\usepackage{longtable}
%% Ragged-right fixed-width column: the table columns here are narrow enough
%% that justification opens large interword gaps.
\newcolumntype{L}[1]{>{\raggedright\arraybackslash}p{#1}}

\newcommand{\EXHALE}{\textsc{Exhale}}
\newcommand{\Kompot}{\textsc{Kompot}}
\newcommand{\Mdot}{\dot{M}}
\newcommand{\Lya}{Ly$\alpha$}
\newcommand{\Ha}{H$\alpha$}
\newcommand{\Kzz}{K_{zz}}
\newcommand{\ergs}{erg\,s$^{-1}$}
\newcommand{\gs}{g\,s$^{-1}$}
\newcommand{\cmc}{cm$^{-3}$}
%% Compact table body: the class sets \linespread{1.5} for running text, which
%% makes multi-line p-columns unnecessarily tall.
\newcommand{\tabsetup}{\footnotesize\linespread{1.05}\selectfont
  \renewcommand{\arraystretch}{1.15}}

\title{The \Kompot\ Code (Johnstone et al.\ 2018, 2019, 2021;\\
Robeling et al.\ 2026) versus \EXHALE:\\
model comparison and implementation-path analysis}
\author{K.-I. Seon}
\date{\docmoddate}

\begin{document}
\maketitle

\section{Purpose and summary}

This memo compares the \Kompot\ Code, as documented across its four defining papers, with
\EXHALE, and then assesses two possible construction paths: writing a \Kompot-class code from
scratch using the published papers, and extending \EXHALE\ with selected \Kompot\ physics. The
headline verdicts, stated first, are the following.

\begin{enumerate}
  \item \textbf{The two codes appear largely complementary rather than competing.} \Kompot\
    targets the thermal and chemical structure of weakly to moderately irradiated thermospheres
    --- terrestrial secondary atmospheres, and in the 2026 paper a hydrostatic H$_2$ giant ---
    while \EXHALE\ targets strongly irradiated transonic H/He\,$+$\,metal winds and the
    spectral line diagnostics that trace them. The two overlap only in the transonic corner,
    which is the subject of Johnstone et al.\ (2019) and the core regime of \EXHALE.

  \item \textbf{The deepest architectural difference is the treatment of chemistry.} \Kompot\
    integrates a time-dependent kinetic reaction network (Rosenbrock solver) and transports every
    species by advection and diffusion; \EXHALE\ solves local ionization equilibrium in each cell
    at every step. \EXHALE's choice appears valid where the ionization and recombination
    timescales are short compared with the flow time, which is its own regime, but it cannot by
    construction represent transport-limited neutral photochemistry such as hydrocarbon or
    CO$_2$ chemistry. That limitation is presumably why \EXHALE\ couples to \textsc{Vulcan}
    offline for the lower atmosphere rather than solving that chemistry internally.

  \item \textbf{From-scratch feasibility differs sharply between the papers.} Johnstone et al.\
    (2018) is unusually complete: all governing equations, numerical appendices, and a reaction
    table running to roughly ten pages. A reimplementation in that operating mode appears
    feasible provided the external data sets are collected. The new numerics introduced in
    Johnstone et al.\ (2019) are underspecified but can be replaced by standard methods. Robeling
    et al.\ (2026), by contrast, does not appear reimplementable from the paper alone: the
    chemical network is not tabulated, the CH$_4$ cooling fit is unpublished, the H$_3^+$
    non-LTE factor is not stated, and the infrared heating channel is never described.

  \item \textbf{For \EXHALE, a few ports look realistic and one does not.} Thermal conduction,
    the closed-form CO$_2$/NO/O cooling functions, and a \Kompot-style subsonic (Jeans-boundary)
    solution mode are all plausible additions. Adopting kinetic chemistry with species transport
    would, in contrast, touch the solver core deeply enough that it would amount to writing a new
    code rather than extending the existing one.
\end{enumerate}

\section{The \Kompot\ lineage and its three operating modes}

The four papers do not describe one fixed code. They describe a common physics core that is run
in three distinct hydrodynamic modes: a semi-static mode with a moving exobase (2018, 2021), a
fully transonic Godunov mode (2019), and a hydrostatic mode (2026). Table~\ref{tab:lineage}
summarizes the lineage; the paragraphs that follow give the detail.

\begin{table}[htbp]
\centering
\caption{The four \Kompot\ papers, their operating modes, and their targets.}
\label{tab:lineage}
\vspace{2pt}
{\tabsetup
\begin{tabular}{@{}L{0.148\textwidth}L{0.148\textwidth}L{0.188\textwidth}L{0.148\textwidth}L{0.253\textwidth}@{}}
\toprule
Reference & Operating mode & Composition and target & Escape regime & Key result \\
\midrule
Johnstone et al.\ (2018), A\&A 617, A107 &
Semi-static, three temperatures, moving exobase &
Full N/O/C/H network, 63 species; modern Earth and Venus &
Jeans escape at the exobase &
Foundation paper; validated against NRLMSISE-00, IRI-2007 and Hedin (1983); solar XUV evolution
from $-3$ to $+2.5$~Gyr \\[3pt]
Johnstone et al.\ (2019), A\&A 624, L10 &
Transonic, MUSCL-Hancock, single temperature &
Earth composition (N$_2$/O$_2$) at 1~AU &
Fully transonic hydrodynamic wind &
Earth's atmosphere under the young active Sun escapes hydrodynamically,
$\Mdot = 1.8\times10^{9}$~\gs \\[3pt]
Johnstone et al.\ (2021), EPSL 576, 117197 &
Semi-static (2018 mode restored) &
Reduced C/N/O network, 30 species; Archean Earth CO$_2$/N$_2$ &
Jeans escape, with a rapid-escape threshold &
At least ${\sim}40\%$ CO$_2$ needed at 3.8~Ga; the young Sun was likely a slow rotator \\[3pt]
Robeling et al.\ (2026), arXiv:2608.04086 &
Hydrostatic ($v = 0$ imposed) &
H$_2$-rich with hydrocarbons to C$_2$H$_6$; Jupiter &
None (bound atmosphere) &
Reproduces the Galileo/JUNO thermal inversion; Joule heating dominates the energy budget as
prescribed \\
\bottomrule
\end{tabular}}
\end{table}

\subsection{Johnstone et al.\ (2018): the foundation}

The 2018 paper defines the physics core. It is a one-dimensional, time-dependent model carrying
three separate temperatures --- neutral, ion, and electron --- coupled by collisional energy
exchange. The hydrodynamics follow the semi-static reduction of Tian et al.\ (2008): mass and
momentum are assumed steady, the velocity ordinary differential equation is integrated downward
from the exobase given an imposed exobase outflow speed taken from the Jeans-consistent
expressions of Luger et al.\ (2015), and the density equation is integrated upward from a fixed
lower boundary. The upper boundary is the exobase itself and moves with the solution; it is
placed where the mean free path equals the pressure scale height, using a constant collision
cross section of $2\times10^{-15}$~cm$^{2}$.

The chemistry is a kinetic network of 63 species, 30 of them ions, connected by 503 reactions
including 56 photoreactions and 7 photoelectron reactions. It is integrated with the RODAS3
Rosenbrock method using an analytic Jacobian, and the chemistry and radiation updates are applied
every 100 hydrodynamic steps rather than every step. The model was validated against NRLMSISE-00
and IRI-2007 for the modern Earth and against Hedin et al.\ (1983) for Venus. Two applications
follow: an experiment varying the CO$_2$ abundance from $0.01\times$ to $1000\times$ the present
value, and an evolutionary sequence in solar XUV from 3~Gyr in the past to 2.5~Gyr into the
future. In the latter, the exobase neutral temperature at 3~Gyr ago is about 4600~K (with the
electron temperature about 5100~K), against a present-day value near 1100~K read from their
figure.

The design principle worth flagging is that \emph{no heating efficiency parameter appears
anywhere}. The heating rate is instead built up term by term: excess energy from photodissociation,
a photoelectron degradation spectrum following Schunk \& Nagy (1978) with 100 logarithmic bins
between 1 and 1000~eV and local deposition, tabulated reaction enthalpies, infrared absorption by
CO$_2$ and H$_2$O using HITRAN/HITEMP line data processed through \texttt{kspectrum}, and a Joule
term calibrated to a global total of $1.4\times10^{18}$~\ergs\ for the Earth with a uniform
$B = 0.5$~G.

\subsection{Johnstone et al.\ (2019): the transonic mode}

The 2019 letter keeps the physics and replaces the hydrodynamics. The semi-static reduction gives
way to a MUSCL-Hancock scheme on primitive variables following Toro (1999), and the three
temperatures collapse to one because the chosen Riemann solver requires a single temperature. The
grid is fixed, spanning altitudes from 50 to 50\,000~km in 500 cells, with a zero-gradient
outflow condition at the top. Notably, no sonic-point condition is imposed anywhere: the
transonic solution emerges from time relaxation.

One model is presented. An Earth-composition (N$_2$/O$_2$) atmosphere is placed at 1~AU under the
young active Sun, using the Claire et al.\ (2012) spectrum for 4.4~Gyr ago, in which the flux
below 100~nm is enhanced by roughly a factor of 60. The result is a fully transonic wind with
$\Mdot = 1.8\times10^{9}$~\gs, which would remove the Earth's present atmosphere in about
0.1~Myr. At the top of the domain the temperature reaches 19\,000~K and the outflow speed 7~km\,s$^{-1}$,
against an escape velocity of 4~km\,s$^{-1}$. The outflow is dominated by atomic N and O and reaches
up to 20\% ionization, and hydrogen is lost about 2.5 times more efficiently than most other
species.

\subsection{Johnstone et al.\ (2021): the Archean application}

The 2021 paper returns to the 2018 semi-static three-temperature mode; the 2019 numerics are not
used. The chemistry is reduced to the C, N and O elements, leaving 30 species with 11 ions and
241 reactions including 42 photoreactions. The lower boundary is a pure CO$_2$\,$+$\,N$_2$
mixture at $T = 267$~K and $n = 5\times10^{16}$~\cmc\ at 50~km altitude. The parameter study
crosses base CO$_2$ mixing ratios from 0.001 to 0.99 with solar activity levels
$F_X \approx 2$--15~\ergs\,cm$^{-2}$ at 1~AU (roughly 3--23 times the modern mean) taken along the
rotation tracks of Tu et al.\ (2015).

Escape is counted as Jeans escape at the exobase only. A ``rapid escape'' regime is diagnosed from
a runaway in the exobase solution together with a mass-loss threshold
$\Mdot > 0.1$~bar\,Myr$^{-1}$, about $1.7\times10^{7}$~\gs. The principal conclusions are that at
3.8~Ga the atmosphere required at least about 40\% CO$_2$ to survive, that above
$F_X \approx 13$~\ergs\,cm$^{-2}$ even a pure CO$_2$ atmosphere escapes rapidly, and that the
young Sun must therefore have been a slow rotator.

\subsection{Robeling et al.\ (2026): the giant-planet extension}

The 2026 paper extends the code to a giant planet and drops the hydrodynamics entirely: $v = 0$ is
imposed in the momentum relation, so the model is hydrostatic, and the exobase is located where
the Knudsen number reaches unity. The target is an H$_2$-rich Jupiter benchmarked against Galileo
and JUNO thermal structure and hydrocarbon profiles. The chemistry is extended to hydrocarbons up
to C$_2$H$_6$, drawn from KIDA, PHIDRATES and several further networks, but the resulting network
is not tabulated in the paper. H$_3^+$ cooling is taken from Table~5 of Miller et al.\ (2013)
multiplied by a non-LTE scale factor whose value is not stated; CH$_4$ cooling comes from an LTE
fit to ExoMol data over 1--1500~K, with the fit itself deferred to a paper in preparation.

Joule heating dominates the model energy budget and is a scaled free parameter. The total
$\Gamma_j \approx 10^{21}$~\ergs\ is taken from M\"uller-Wodarg et al.\ (2025) and distributed
over a shape function computed from the Pedersen conductivity with a placeholder unit electric
field and a single scalar $B = 4.17$~G. The eddy diffusion coefficient $\Kzz$ is the other major
free parameter, explored using the Moses et al.\ (2005) profiles~A and~C together with constant
values of $10^{6}$ and $10^{7}$~cm$^{2}$\,s$^{-1}$. The model reproduces the thermal inversion and
the upper-atmosphere temperatures well, overpredicts C$_2$H$_2$ by a factor of about 2--3, and the
authors acknowledge that aerosols, particle precipitation and non-LTE CH$_4$ are missing. A
coupling to \textsc{TauREx} produces gas-only transmission spectra over 1--15~$\mu$m, in which
CH$_4$ dominates from 1 to 10~$\mu$m and C$_2$H$_2$ beyond 10~$\mu$m.

Two internal tensions are worth recording, tentatively. First, the headline that Joule heating
dominates appears to follow from the input choice of $\Gamma_j$ rather than from an independent
derivation, since the total is prescribed rather than computed. Second, Table~1 quotes
$\Gamma_j \approx 10^{21}$~\ergs, which is 100~TW, while the discussion of Fig.~3 reports the
integrated Joule contribution as about 200~TW; the two statements do not appear reconcilable from
the text alone.

\section{Physics comparison with \EXHALE}

The verdict first: within \EXHALE's own regime --- strongly irradiated escaping H/He\,$+$\,metal
atmospheres and the line diagnostics that probe them --- nothing in the four \Kompot\ papers
supersedes what \EXHALE\ already does, and in several places (\Lya\ transfer, excited hydrogen,
metal line cooling, transit spectra) \EXHALE\ goes considerably further. Conversely, in \Kompot's
regime --- bound or weakly escaping thermospheres with neutral molecular chemistry --- \EXHALE\
currently cannot operate at all, because it has no bound-atmosphere solution mode, no thermal
conduction, and no kinetic chemistry. Table~\ref{tab:physics} sets the two side by side. The
\EXHALE\ entries were checked against the source tree rather than against its documentation.

{\tabsetup
\begin{longtable}{@{}L{0.155\textwidth}L{0.39\textwidth}L{0.39\textwidth}@{}}
\caption{Physics comparison. \Kompot\ entries are as described in the four papers; \EXHALE\
entries were verified in \texttt{EXHALE/src/}.}
\label{tab:physics}\\
\toprule
Aspect & \Kompot\ (2018--2026) & \EXHALE \\
\midrule
\endfirsthead
\multicolumn{3}{@{}l}{\emph{Table \ref{tab:physics}, continued}}\\
\toprule
Aspect & \Kompot\ (2018--2026) & \EXHALE \\
\midrule
\endhead
\bottomrule
\endfoot
Hydrodynamics and regime &
Three modes: semi-static with moving exobase (2018, 2021), transonic MUSCL-Hancock (2019),
hydrostatic $v=0$ (2026). Bound and escaping atmospheres both representable. &
Single-fluid three-equation Euler system, spherical, Roche potential on by default. Transonic
wind only; a bound hydrostatic solution regime does not exist in the code. \\

Steady-state method &
Time relaxation in all modes; convergence criterion not stated in any of the papers. &
Explicit SSP-RK3 marching, optionally finished by a pseudo-transient-continuation
Jacobian-free Newton--Krylov solver with GMRES
(\texttt{steady\_newton.f90}). Convergence on the mass-flux spread
$du = \mathrm{spread}(\rho \varv r^2) < 10^{-3}$. \\

Numerical scheme &
2019: MUSCL-Hancock on primitive variables, Riemann solver cited only by equation number in
Toro (1999). 2018/2021: ODE integration, no Riemann solver. &
PLM or WENO3 reconstruction; HLLC, Roe or LLF fluxes (\texttt{Num\_Fluxes.f90},
\texttt{Reconstruction.f90}). \\

Temperatures &
Three ($T_n$, $T_i$, $T_e$) with collisional exchange in 2018, 2021, 2026; collapsed to one in
2019 because the Riemann solver requires it. &
One temperature shared by neutrals, ions and electrons,
$T = p / [(n_{\rm tot} + n_e)k]$. No $T_e$ anywhere in the source. \\

Thermal conduction &
Present in all modes, with separate conductivities for neutrals, ions and electrons; dominant
near the base of a bound thermosphere. &
Available but off by default. The conduction slot in the hydrodynamic source vector is
identically zero (\texttt{Source.f90}, \texttt{S(3) = 0.0}); conduction and viscosity are
instead an operator-split Crank--Nicolson step
(\texttt{viscous\_conduction.f90}, keys \texttt{Conduction:} and \texttt{Viscosity:},
both default off), applied in the marching loop and evaluated by the same routine in the
steady residual so that the two agree. One conductivity for the single fluid,
$\kappa(T)=4.45\times10^{4}\,(T/1000\,\mathrm{K})^{0.7}$ (Watson et al.\ 1981), with
$\mu(T)$ tied to it by the monatomic Chapman--Enskog relation
(\texttt{docs/viscosity\_conduction.md}). \\

Photo-heating and heating efficiency &
Built from first principles; no heating efficiency parameter anywhere, stated explicitly as a
design goal. &
Photoionization heating computed from the frequency integral, also with no efficiency parameter.
The $\eta$ reported in the output is a diagnostic, not an input. The design stance is the same in
both codes. \\

Secondary and photoelectron treatment &
Photoelectron degradation spectrum transported over 100 logarithmic bins from 1 to 1000~eV
(Schunk \& Nagy 1978), with local deposition. &
Secondary ionization by fast photoelectrons through the Shull \& van Steenberg (1985) fits: a
local energy-partition treatment, not a degradation-spectrum calculation. \\

Chemical heating &
General bookkeeping of tabulated reaction enthalpies across the whole network. &
Only specific channels: He($2^3$S)\,$+$\,H and He($2^3$S)\,$+$\,H$_2$ Penning ionization (6.2 and
4.4~eV), He recombination-radiation coupling, and photoelectrons from H($n=2$) Balmer-continuum
ionization. No general reaction-enthalpy bookkeeping. \\

Joule heating &
Present. Earth: total $1.4\times10^{18}$~\ergs, $B = 0.5$~G. Jupiter: total
$\Gamma_j \approx 10^{21}$~\ergs\ distributed over a Pedersen conductivity shape,
$B = 4.17$~G. &
Absent. A search of the source for Joule, magnetic or Pedersen terms returns no hits; there is
no magnetic field in the model. \\

Infrared and bolometric heating &
CO$_2$ and H$_2$O infrared absorption (2018) from HITRAN/HITEMP; an infrared heating channel is
referred to but not described in 2026. &
No bolometric infrared heating in the production solve. The energy grid floor is 4.8~eV (or the
lowest metal ionization potential); the only sub-XUV stellar inputs are a diluted-blackbody
Balmer continuum acting on H($n=2$) and stellar \Lya\ pumping. \\

Cooling channels &
CO$_2$ 15~$\mu$m cool-to-space (two-level, Kumer \& James escape probabilities), NO 5.3~$\mu$m,
O 63/147~$\mu$m, H$_2$O; plus H$_3^+$ and CH$_4$ in 2026. &
\Lya; H and He recombination, collisional ionization and free-free; He\,$2^3$S channels; metal
lines as closed-form CHIANTI~v11 fits for C\,\textsc{i/ii}, N\,\textsc{i/ii}, O\,\textsc{i/ii},
Mg\,\textsc{i/ii}, Ca\,\textsc{ii}, Na\,\textsc{i}, Fe\,\textsc{i/ii}, with density-dependent
[C\,\textsc{ii}]~158~$\mu$m and [O\,\textsc{i}]~63~$\mu$m saturation floors and a
statistical-equilibrium Fe\,\textsc{ii} coefficient; H$_3^+$ infrared cooling
(\texttt{h3p\_cooling.f90}). No CO$_2$, NO or CH$_4$ cooling. \\

H$_3^+$ cooling detail &
Miller et al.\ (2013) Table~5 fits times a non-LTE scale factor whose value is not given
(2026). &
The same Miller et al.\ (2013) Table~5 fits, plus the Table~6 non-LTE departure factor. On this
one item \EXHALE\ appears better documented than Robeling et al.\ (2026). \\

Chemistry treatment &
Time-dependent kinetic network integrated with RODAS3 Rosenbrock and an analytic Jacobian; every
species carries its own continuity equation. &
Local steady-state ionization equilibrium solved in every cell at every step
(\texttt{util\_ion\_eq.f90}, \texttt{ionization\_equilibrium.f90}), by an analytic-Jacobian
Newton method with a MINPACK \texttt{hybrd1} fallback. \\

Species set &
2018: 63 species, 30 ions, 503 reactions. 2021: 30 species, 11 ions, 241 reactions. 2026: extended
to hydrocarbons up to C$_2$H$_6$ (network not tabulated). &
37 species: H\,\textsc{i/ii}; He\,\textsc{i/ii/iii} and He\,$2^3$S; 27 metal ion stages of the ten
elements C, N, O, Mg, Si, Ca, Na, K, S, Fe (\texttt{species\_table.f90}); and H$_2$, H$_2^+$,
H$_3^+$, HeH$^+$. Charge exchange from Huang et al.\ (2023) Table~4; the molecular network is 23
reactions with Koskinen et al.\ (2022) rates, also in local equilibrium. H$_2$ Lyman--Werner
photodissociation is absent, as noted in the code comments. \\

Species transport &
Every species is advected and diffused; molecular diffusion coefficients from Banks \& Kockarts
with an N$_2$ background, plus eddy diffusion. &
No species advection in the main loop; an advection-corrected re-solve exists only as a
post-processing step (\texttt{post\_process\_adv.f90}). Binary H/He element transport of the
helium MASS FRACTION $X=\rho_{\rm He}/\rho$ with molecular $D_{12}$, eddy $\Kzz$,
ambipolar-corrected settling (on by default within the option), the molecular-carrier
term $-d\ln\psi/dr$ and optional thermal diffusion
(\texttt{binary\_element\_diffusion.f90}); each metal may additionally diffuse as a trace
species against the hydrogen carriers (\texttt{He\_metal\_diffusion}, off by default).
The whole diffusion option is off by default, and there is no eddy diffusion of the
bulk gas. \\

Radiative-transfer geometry &
Single slant ray at a chosen zenith angle (66$^\circ$ for Earth, 60$^\circ$ in 2026); optical
depth summed over all photoreactions. &
Radial column with $\exp(-\tau)$ attenuation and a day--night redistribution factor $\xi$ (0.5 or
0.25). Spectrum from a broken power law, a tabulated SED, or monochromatic; default range
13.6--1240~eV (the floor drops to 4.8~eV with He\,$2^3$S, or to the lowest neutral-metal
ionization potential), in about 200 bins. The two treatments are different one-dimensional
averaging approximations, not a right-versus-wrong distinction. \\

\Lya\ transfer and excited hydrogen &
Not treated. &
\Lya\ radiative transfer through a Neufeld escape-probability closure computed in line, or an
imported mean-intensity profile from the LaRT Monte Carlo code (\texttt{lya\_rt.f90}); full
non-LTE H($n=2$) 2s/2p balance feeding Balmer-continuum ionization and heating
(\texttt{excited\_hydrogen.f90}). Nothing comparable exists in \Kompot. \\

Boundaries and domain &
Lower boundary fixed at 50~km (2019, 2021) or 65~km (2018 Earth). Upper boundary at the moving
exobase (2018, 2021, 2026) or zero-gradient outflow at 50\,000~km (2019). &
500 cells by default from the base radius to the Hill or L1 radius set by the Roche geometry
(overridable). Fixed base density and temperature, with $T_0$ the equilibrium temperature, plus a
one-way valve enforcing $\varv \ge 0$ at the base; zero-gradient outflow at the top; $\Mdot$ read
near the top. \\

Free parameters &
Eddy $\Kzz$, described by the authors as the only real free parameter (2019); the Joule total and
its electric field; the H$_3^+$ non-LTE factor (2026). &
Eddy $\Kzz$ when diffusion is enabled (the scalar \texttt{He\_Kzz}, default 0, or a
$\Kzz(p)$ column of a \texttt{Lower atmosphere profile}), and the day--night redistribution
factor $\xi$. \\

Outputs and observables &
Thermal, chemical and ionization structure; exobase state; Jeans or hydrodynamic mass-loss rate;
gas-only transmission spectra through \textsc{TauREx} in 2026. &
Thermal, chemical and ionization structure; $\Mdot$; and transit transmission spectra for
He\,\textsc{i}~10830, \Lya, \Ha, H$\beta$, Mg\,\textsc{ii}~h\&k, Ca\,\textsc{ii}~H\&K and
Na\,\textsc{i}~D. \\
\end{longtable}}

\subsection{Free parameters, stated honestly}

The \Kompot\ papers make a point of using no heating efficiency parameter, and that claim appears
accurate. It should not be read, however, as meaning that the models are free of tuned inputs. By
the authors' own statement the eddy diffusion coefficient is a free parameter; the Joule heating
total and the electric field used to shape its deposition are prescribed rather than derived; and
the H$_3^+$ non-LTE scale factor in the 2026 paper is an unstated multiplier. \EXHALE\ shares the
first of these ($\Kzz$, when diffusion is switched on) and adds the day--night redistribution
factor $\xi$, which plays a role loosely analogous to \Kompot's choice of zenith angle. On the
specific question of heating efficiency the two codes take the same first-principles position.

\section{Reproducibility of each paper}

Verdict first: Johnstone et al.\ (2018) comes close to a complete specification of a working code;
the 2019 and 2021 papers inherit that specification but leave their own new pieces
underdocumented; and Robeling et al.\ (2026) does not appear reproducible from the paper alone.

\subsection{Johnstone et al.\ (2018)}

What is fully specified is unusually broad. It includes all governing equations, including the
semi-static reduction and the directions in which each ODE is integrated; the radiation scheme
(10--4000~\AA\ in 1000 bins, infrared 1--20~$\mu$m in roughly $10^4$ selected bins, ray step
$H/5$, optical depth summed over photoreactions); the photoelectron degradation equation with its
bin structure and all cross-section sources; every heating and cooling formula together with its
constants (the CO$_2$ cool-to-space two-level model with Kumer \& James escape probabilities and
the collider rate table, NO 5.3~$\mu$m, O 63 and 147~$\mu$m, and the conductivities for neutrals,
ions and electrons); the diffusion coefficients and thermal-diffusion factors; all numerical
solvers at appendix level of detail (Crank--Nicolson tridiagonal forms, RODAS3 with tolerances,
MUSCL/TVD parameters); the complete reaction table with rates, enthalpies and validity ranges
(Table~H.1); the operator-splitting order and the 100-step update cadence; and a set of solver
validations against an analytic Parker wind, a Sod tube compared with VAC, a cross-check against
KROME, and a photoelectron spectrum compared with Singhal \& Haider (1984).

What is left out is a shorter but not empty list: the number of grid cells and the exact spacing
law; the steady-state convergence criterion, which is likely the largest single obstacle to
reproducing the quoted numbers; the Courant number, given only as $C \sim 1$; the Joule electric
field, which is tuned to a global total and is therefore not reconstructible for a single column;
the H$_2$O cooling equations, for which the reader is pointed to Kasting \& Pollack (1983); the
algorithm for the thermal/non-thermal crossover energy $E_t$; the energy-loss tabulations
$\Delta E$; the ion and electron cooling terms $Q_{c,i}$ and $Q_{c,e}$, which appear in the
equations but are never given; and the exobase regridding algorithm. Beyond the paper, an
implementer would need the external numerical arrays: PHIDRATES cross sections; electron-impact
cross sections from Green \& Barth (1965), Watson et al.\ (1967), Jackman et al.\ (1977) and
Jusick et al.\ (1967); the Banks \& Kockarts and Schunk \& Nagy collision tables; HITRAN/HITEMP
line lists processed through \texttt{kspectrum}; and the Claire et al.\ (2012) spectra.

\subsection{Johnstone et al.\ (2019)}

The physics is inherited from 2018 and is therefore documented. What the letter adds is a change
of numerics, and that change is underspecified: the Riemann solver is identified only by an
equation number in Toro (1999); the limiter, the operator-splitting order, the CFL number, the
sonic-point cutoff used to disable diffusion, the definition used for the mass-loss rate, and the
convergence criterion are all absent. The 50~km base state --- temperature, density, and all 63
mixing ratios --- is not tabulated, and it cannot simply be copied from the 2018 paper because
that Earth model was based at 65~km. None of these gaps involves physics, so all of them can be
closed by an implementer's own reasonable choices, at the cost of not reproducing the quoted
numbers exactly.

\subsection{Johnstone et al.\ (2021)}

This paper is largely derivable from 2018 by mechanical means: the element filter rule is given,
and the CO$_2$/N$_2$ base state is in fact better documented than the 2019 base state. The real
gaps are the eddy diffusion coefficient adopted for a CO$_2$-rich atmosphere, which directly
controls how much CO$_2$ reaches the radiating region and therefore controls the headline
threshold; whether Joule heating was enabled; the grid; the set of species allowed to escape by
Jeans escape; and the conversion used between bar\,Myr$^{-1}$ and \gs.

\subsection{Robeling et al.\ (2026)}

Four items block an independent implementation, and for each of them there is no citable closing
source. The chemical network is not tabulated, and it is described as the union of six cited
sources with an undocumented selection rule. The CH$_4$ cooling fit is deferred to a paper in
preparation. The H$_3^+$ non-LTE scale factor is given neither a value nor a reference. The
infrared heating channel appears in the abstract and in a figure legend but is never described.
Secondary omissions are the epoch of the solar spectrum and the orbital distance, and the internal
100~TW versus 200~TW inconsistency noted in \S2.4. The conclusion is that an independent
implementation could only be built ``in the style of'' Robeling et al.\ (2026) --- for instance,
with a \textsc{Vulcan}-derived network and fits of the implementer's own construction --- and
should not be described as a reproduction.

\section{Implementation paths}

\subsection{Building a \Kompot-class code from scratch}

A reimplementation in the 2018 operating mode appears feasible as a standalone project. The paper
together with its appendices amounts to a near-complete specification, and the missing numerical
items (grid, convergence criterion, Courant number) are choices any implementer can make without
guessing at physics. The labor would be dominated not by the equations but by collecting the
external data --- cross-section and rate archives, collision tables, line lists --- and by
entering the 503-reaction network. Reproducing the published curves qualitatively appears
realistic on that basis; matching the quoted numbers exactly does not, because the convergence
criteria are not published. Adding the transonic mode of 2019 requires an independent choice of
Godunov-type numerics, for which \EXHALE-grade methods would be entirely adequate. None of this
requires inventing new physics.

\subsection{Extending \EXHALE}

Table~\ref{tab:ports} lists the candidate ports with a verdict, a rough effort estimate, and an
assessment of physical value; the notes below expand on the entries that need it.

\begin{table}[htbp]
\centering
\caption{Candidate \Kompot-inspired additions to \EXHALE.}
\label{tab:ports}
\vspace{2pt}
{\tabsetup
\begin{tabular}{@{}L{0.03\textwidth}L{0.235\textwidth}L{0.10\textwidth}L{0.55\textwidth}@{}}
\toprule
& Candidate & Effort & Verdict and physical value \\
\midrule
(a) & CO$_2$/NO/O closed-form cooling (2018 Eqs.~34--43 plus the collider table) & low &
Trivial to code, but meaningless without the corresponding neutral species. Valuable only together
with~(d). \\
(b) & Thermal conduction & moderate &
Realistic port. All coefficients are in the 2018 paper; for a single-temperature gas the
density-weighted mixing of 2019 applies. Prerequisite for~(c). \\
(c) & Subsonic/Jeans solution mode & large &
The most scientifically interesting port. Would extend \EXHALE\ to weakly irradiated planets where
no transonic wind exists. \\
(d) & Kinetic chemistry with species transport & very large &
Not an extension. Would effectively be a new code. \\
(e) & Schunk \& Nagy photoelectron degradation & moderate &
Closed specification, but marginal gain in \EXHALE's strongly ionized regime. \\
(f) & Joule heating in the Robeling form & low &
Technically easy, but imports an unconstrained free parameter. Sensitivity studies only. \\
(g) & Robeling-style Jupiter infrared thermosphere & --- &
Out of reach as \EXHALE\ stands; needs (d) and (b) plus a hydrostatic mode. \\
\bottomrule
\end{tabular}}
\end{table}

\paragraph{(b) Thermal conduction.}
The effort is moderate rather than small because the term must enter in two places: \EXHALE\
already has an implicit energy-update pattern that can be extended, but the Jacobian-free
Newton--Krylov residual would also need the conduction term for the steady solver to remain
consistent. Physically, conduction is negligible where advection dominates, which is the case in
\EXHALE's converged winds, but it matters near the base and it matters throughout for weakly
irradiated planets. That is why it is a prerequisite for~(c) rather than an end in itself.

\paragraph{(c) A subsonic/Jeans solution mode.}
This means a semi-static hydrodynamic solve with a moving exobase and a Jeans-consistent upper
boundary, and it is structurally the largest realistic addition on the list. The algorithm is
documented completely in Appendix~D of the 2018 paper, so the specification risk is low even though
the implementation effort is high. The scientific argument for it is that \EXHALE\ currently
cannot represent a bound thermosphere at all: the base valve and the outflow top together assume a
wind. Adding this mode would extend the code's planet class downward, to planets where a transonic
wind does not exist.

\paragraph{(d) Kinetic chemistry with species transport.}
This is the chasm in the comparison. Replacing local ionization equilibrium with a Rosenbrock
network integration plus advection and diffusion of every species would touch the solver core, the
Newton--Krylov residual, and the post-processing chain simultaneously. It is better described as
writing a new code than as extending \EXHALE. If terrestrial-composition escape --- Archean-Earth
type questions --- ever becomes a research goal here, a fresh implementation guided by the 2018
paper appears more sensible than bending \EXHALE\ to fit.

\paragraph{(f) and (g).}
Joule heating is easy to add in the Robeling form, since it amounts to distributing a prescribed
total over a Pedersen conductivity shape, but for exoplanets both $B$ and the effective electric
field are unknown, so its value appears limited to sensitivity studies. As for a Robeling-style
Jupiter application, it is out of reach for \EXHALE\ as it stands. It is worth noting, though, that
\EXHALE\ already has the H$_3^+$ cooling piece, and that it reaches the neutral molecular lower
atmosphere by a different route --- the offline \textsc{Vulcan} coupling --- rather than by
solving that chemistry internally.

\subsection{Recommendation}

Tentatively, the comparison suggests keeping the two codes' regimes separate rather than trying to
make either one cover both. For \EXHALE's own science, (b) and (c) appear to be the natural
\Kompot-inspired upgrades: they extend the range of planets the code can address without
disturbing what it already does well. Items (a) and (d) are worth revisiting only if secondary
atmospheres become a research goal, and in that case a new code guided by the 2018 paper looks
more sensible than an extension; item (f) is worth having only as a sensitivity knob. Finally,
the \Kompot\ papers independently validate \EXHALE's design stance on first-principles heating:
both codes compute heating term by term, neither uses a heating efficiency parameter, and the
\Kompot\ authors treat that as a selling point.

\section{Caveats}

This memo rests on the four papers --- published text for the three Johnstone papers, preprint
text for Robeling et al.\ (2026) --- together with the \EXHALE\ source tree. The \Kompot\ source
is not public, so every statement here about \Kompot\ is a statement about what its papers say,
and the code may contain physics or numerics they do not describe. Several numbers above are read
from published figures rather than tabulated values: some Johnstone et al.\ (2021) thresholds,
the present-day exobase temperature in Johnstone et al.\ (2018), and the Robeling et al.\ (2026)
temperatures. The effort estimates in \S5.2 are qualitative. The \EXHALE\ statements were checked
in the source at the time of writing and describe the default configuration; several features
discussed are opt-in.

\end{document}
